\chapter{Pure preimages and common finite observations}
\label{ch:purity}
\chapterintro{The common-completion phenomenon is not confined to one selector construction. It follows from a general correspondence between rational subspaces and pure intermediate groups.}

\section{The preimage construction}
Keep $D\subseteq P$ and $V=P/D$ from Chapter~\ref{ch:ambient}. For a rational vector subspace $W\le V$, define
\[
H_W=q^{-1}(W).
\]
This is an actual subgroup of $P$ containing $D$. Its quotient by $D$ is canonically $W$. It is dense because it contains the dense subgroup $D$.

The construction separates the ambient coordinate space from the admissible realization. The quotient $V$ records infinite coordinate tails modulo finite changes. A rational subspace $W$ selects which such tails are allowed. The preimage restores all finite-support corrections.

\begin{definition}[Purity]
A subgroup $H$ of an abelian group $M$ is pure when
\[
H\cap mM=mH\qquad(m\ge1).
\]
Equivalently, if an equation $my=h$ with $h\in H$ has a solution in $M$, then it has a solution in $H$.
\end{definition}
Purity is a standard algebraic notion \cite{Fuchs}. The proof below uses only the displayed criterion. In particular, no structure theorem for mixed abelian groups is required.

\begin{theorem}[Pure-preimage correspondence]
\label{thm:purecorrespondence}
The map
\[
W\longmapsto H_W=q^{-1}(W)
\]
is an inclusion-preserving bijection between rational subspaces of $V$ and pure subgroups $H$ of $P$ containing $D$.
\end{theorem}
\begin{proof}
First let $W$ be rational and suppose $h\in H_W\cap mP$. Because $W$ is closed under division by $m$, choose $x\in H_W$ with $q(x)=q(h)/m$. Then $d=h-mx$ belongs to $D$. At a prime dividing $m$, the coordinate of $h$ is zero because $h\in mP$, and that of $mx$ is also zero. Thus $d_p=0$ whenever $p\mid m$.

Define $e_p=m^{-1}d_p$ at every other prime and $e_p=0$ at primes dividing $m$. The vector $e$ has finite support and satisfies $me=d$. Hence $h=m(x+e)$ with $x+e\in H_W$. This proves purity.

Conversely, let $D\subseteq H\subseteq P$ be pure. Its image $q(H)$ is an additive subgroup of $V$. Given $h\in H$ and $m\ge1$, let $d\in D$ agree with $h$ on the primes dividing $m$ and vanish elsewhere. Then $h-d\in mP\cap H$. Purity gives $h-d=mx$ for some $x\in H$, so $q(x)=q(h)/m$. Thus $q(H)$ is closed under rational scalars. Since $D\subseteq H$, the equality $H=q^{-1}(q(H))$ holds. This proves that the constructions are inverse. Both preserve inclusion.
\end{proof}
The bottom object is $H_{\{0\}}=D$, and the top object is $H_V=P$. Intersections of rational subspaces correspond to intersections of their preimages. Sums correspond to subgroup sums, with $D$ serving as the bottom of the interval of intermediate groups.

\section{Root conditions do not depend on the subspace}
\begin{theorem}[Ambient root criterion]
\label{thm:ambientroot}
For every rational subspace $W$, every $h\in H_W$, and every $m\ge1$,
\[
h\in mH_W\quad\Longleftrightarrow\quad
h_p=0\text{ for all primes }p\mid m.
\]
\end{theorem}
\begin{proof}
Multiplication by $m$ on $P$ is zero at primes dividing $m$ and invertible on every other coordinate. Thus the right-hand condition is equivalent to $h\in mP$. Apply purity from Theorem~\ref{thm:purecorrespondence}.
\end{proof}
The condition is independent of $W$, but the set of admissible $h$ is not. This distinction is the reason identical quotient laws can coexist with different original groups.

\begin{theorem}[One finite-quotient system for every preimage]
\label{thm:generalquotients}
For every rational subspace $W$ and every $m\ge1$, the coordinate map induces an isomorphism
\[
H_W/mH_W\ \cong\ \prod_{p\mid m}C_p.
\]
For $m\mid n$, the quotient transition map is the coordinate restriction from primes dividing $n$ to primes dividing $m$. These quotients are cofinal among all finite quotients, and
\[
\widehat{H_W}\cong P.
\]
\end{theorem}
\begin{proof}
The coordinate map onto the displayed finite product is surjective already on $D$. Its kernel is $mH_W$ by Theorem~\ref{thm:ambientroot}. The transition assertion follows by evaluating the same coordinates.

If $H_W\to F$ is a homomorphism to a finite group, then the image is abelian and is killed by $\lvert F\rvert$. Its kernel therefore contains $\lvert F\rvert H_W$. Hence the displayed kernels are cofinal. Their inverse limit is the product of all prime coordinates, which is $P$. Under this identification the completion map is the original inclusion $H_W\hookrightarrow P$.
\end{proof}
There is a stronger natural statement. It compares maps to every finite target rather than merely comparing the quotient groups.

\begin{proposition}[Finite-target equivalence]
\label{thm:finitehom}
For every finite group $F$, restriction along $D\hookrightarrow H_W$ gives a natural bijection
\[
\Hom(H_W,F)\longrightarrow\Hom(D,F).
\]
\end{proposition}
\begin{proof}
Put $m=\lvert F\rvert$. Every $h\in H_W$ agrees modulo $mH_W$ with an element of $D$, namely its finite coordinate vector on primes dividing $m$. Every homomorphism to $F$ kills $mH_W$. Thus two maps agreeing on $D$ agree everywhere.

Conversely, a map $D\to F$ kills every prime factor whose order does not divide $m$. It factors through the finite coordinate product on primes dividing $m$. Composing that finite-coordinate rule with the coordinate projection of $H_W$ extends it. The extension is unique by the first part. Postcomposition with a map of finite groups commutes with restriction, giving naturality.
\end{proof}
The argument never subtracts maps to a nonabelian group. The assertion is a bijection of sets of homomorphisms, and naturality is with respect to postcomposition.

\section{Countability, rank, and nonsplitting}
\begin{proposition}[Structure of an ambient preimage]
\label{thm:HWstructure}
For every rational subspace $W$,
\[
\Tor(H_W)=D,\qquad H_W/D\cong W.
\]
The group $H_W$ is reduced and residually finite. It is countable exactly when $W$ has finite or countably infinite dimension over $\Q$. If $W\ne0$, the sequence
\[
0\longrightarrow D\longrightarrow H_W\longrightarrow W\longrightarrow0
\]
does not split.
\end{proposition}
\begin{proof}
The torsion equality follows from $D\subseteq H_W\subseteq P$ and Theorem~\ref{thm:torP}. The quotient assertion is the first isomorphism theorem. Reducedness follows from reducedness of $P$. A nonzero element has a nonzero finite coordinate, which detects it and proves residual finiteness.

A rational vector space on a finite or countable basis is countable, since each vector has a finite rational linear expression. Each fiber of $H_W\to W$ is a coset of the countable group $D$. Thus $H_W$ is countable in this case. Conversely, a quotient of a countable group is countable and has at most a countable basis.

A splitting would embed the nonzero divisible group $W$ into $P$. Theorem~\ref{thm:Vrational} excludes every such subgroup.
\end{proof}
The rank of the torsion-free quotient is exactly $\dim_{\Q}W$. This gives countable examples of every finite rank and of countably infinite rank with the same finite-quotient system. The rank-one family studied next is chosen because it admits an especially sharp classification and computability theorem.

\section{What has been fixed and what still varies}
The entire product $P$, all finite quotient objects, their transition maps, and the maps to arbitrary finite targets are fixed throughout the construction. The subspace $W$ remains free. This is an ambient realization theorem, not a method for discovering $W$ from those observations.

The preimage formula also identifies the correct role of completion. Every $H_W$ contains the same dense torsion subgroup. Its completion admits every compatible coordinate tuple. Its original image admits only those tuples whose tail class lies in $W$. The missing information is an admissibility condition on the quotient, not a defect in the finite coordinates themselves.
